\documentclass[11pt]{article}
\usepackage[margin=1in]{geometry}
\usepackage{amsmath,amssymb,amsthm,booktabs,hyperref,listings}
\hypersetup{colorlinks=true,urlcolor=blue,linkcolor=blue}
\newtheorem{theorem}{Theorem}
\newtheorem{lemma}{Lemma}
\newcommand{\Geo}{\operatorname{Geo}}
\title{Three spectral invariants do not determine\newline geodesic growth of tree-based RACGs}
\author{Research draft for MathDB 330264}
\date{31 October 2026 campaign; round 2 dossier prepared 4 November 2026 UTC}
\begin{document}
\maketitle
\begin{abstract}
We give two explicit trees on 19 vertices with equal adjacency, complement and line-graph characteristic polynomials whose right-angled Coxeter groups have different geodesic growth with respect to their vertex generators. At length 18 the numbers of individual shortest words are respectively
\(37769749324097905710910\) and \(37769749324097905710906\).
A complete local Lean certificate proves the negation of the universal finite-tree implication in Ciobanu--Kolpakov, arXiv:1504.02774v3, Conjecture 1.3. The formal chain includes the presented-group normal form, automaton recognition, and a bridge to cardinalities of individual words. This draft makes no minimality, infinite-family, new-normal-form or global-priority claim. The post-publication arXiv revision retains the conjecture; the 2016 journal body remains unread, a qualified source-version limitation.
\end{abstract}
\section{Statement and scope}
For a finite simple graph \(T\) on \(V\), put
\[
 W_T=\langle s_v\ (v\in V)\mid s_v^2=1,\quad
 s_us_v=s_vs_u\ \text{if }\{u,v\}\in E(T)\rangle.
\]
Let \(g_T(k)\) count words of length \(k\) in these generators that have minimum length among all words representing the same element of \(W_T\). Different words representing the same element are counted separately. Let \(\chi_T(x)=\det(xI-A_T)\); write \(\overline T\) for the simple complement and \(L(T)\) for the graph on unordered edges, adjacent when they share a vertex.
\begin{theorem}\label{thm}
There are trees \(G,H\) on \(\{0,\ldots,18\}\) such that
\[
\chi_G=\chi_H,\qquad \chi_{\overline G}=\chi_{\overline H},\qquad
\chi_{L(G)}=\chi_{L(H)},\qquad g_G(18)-g_H(18)=4.
\]
In particular the three spectral invariants jointly fail to determine the full geodesic-growth series even among trees.
\end{theorem}
Theorem~\ref{thm} refutes the universal implication in the inspected arXiv source [1, Conjecture 1.3], as recorded in MathDB 330264. A single unequal coefficient is enough; no asymptotic growth-rate assertion is involved. The journal bibliographic identity is J.~Algebraic Combinatorics 44(2), 249--264 (2016), DOI 10.1007/s10801-016-0667-9. The 2020 arXiv revision is not a 2020 journal publication. The claim concerns the authenticated arXiv conjecture. The author-submitted v3 (17 January 2020), after the 2016 publication, retains its hypotheses and conclusion. Root has resolved the source hold as a qualified limitation; the version of record body remains unread.
\newpage
\section{Explicit witness and spectral certificates}
Use the following unordered edge lists, in the displayed order:
\begin{align*}
E(G)=\{&01,12,19,0\!:\!10,23,34,45,56,58,67,\\
       &10\!:\!11,11\!:\!12,11\!:\!17,12\!:\!13,13\!:\!14,14\!:\!15,15\!:\!16,17\!:\!18\},\\
E(H)=\{&01,12,0\!:\!11,23,29,34,45,56,58,67,9\!:\!10,\\
       &11\!:\!12,12\!:\!13,12\!:\!18,13\!:\!14,14\!:\!15,15\!:\!16,16\!:\!17\}.
\end{align*}
Here \(ij\) denotes the edge \(\{i,j\}\) for single-digit labels; a colon separates multi-digit labels. Machine-readable ordered pairs are supplied in the dossier:
\begin{quote}\small\ttfamily sources/spectral\_candidate.json\\
proof/SpectralDefinitions.lean\end{quote}
Both graphs have 18 edges. Paths from 0 reach every vertex (follow the branches in the lists); hence they are connected trees.

For explicit arithmetic verification, the common characteristic-polynomial coefficient lists below are in descending degree, including zero coefficients. They specify degrees 19, 19 and 18 respectively.
\begin{quote}\small\ttfamily
Adjacency:\newline
[1,0,-18,0,133,0,-522,0,1179,0,-1558,0,1177,0,\newline
-473,0,85,0,-4,0]\par
Complement:\newline
[1,0,-153,-1366,-5283,-9226,-255,26518,36212,\newline
-5482,-49704,-28056,19616,22666,847,-5160,\newline
-1273,156,42,0]\par
Line graph:\newline
[1,0,-20,-6,161,88,-658,-484,1433,1242,-1637,\newline
-1542,937,892,-249,-212,24,14,-1]
\end{quote}
These lists can be checked entirely in rational arithmetic. For any of the six adjacency matrices \(A\), set \(B_0=I\), and for \(1\leq k\leq n\) compute
\[
 C_k=AB_{k-1},\qquad c_k=-\operatorname{tr}(C_k)/k,\qquad
 B_k=C_k+c_k I.
\]
The Newton identities give \(\det(xI-A)=x^n+c_1x^{n-1}+\cdots+c_n\).
The dossier includes the exact arithmetic script and coefficient JSON; no numerical eigenvalue comparison is used.

The Lean proof uses a second certificate form: explicit rational matrices \(P,Q\) with \(PQ=QP=I\) and \(AP=PB\) for each pair of matrices. Then \(A=PBQ\), and similarity preserves characteristic polynomials. The split spectral modules supply the entries and row certificates, assembled by \texttt{SpectralWitness.lean}. For line graphs, the edge-list enumeration is proved bijective onto the actual unordered edges; preservation of shared-endpoint adjacency gives an isomorphism with Mathlib's intrinsic line graph. Thus the third identity concerns actual line graphs. The ordinary coefficient computation is a corroborating route, not an assumption in that certificate.
\newpage
\section{Shortest words in the actual group}
Call a word \emph{clean} when it has no two occurrences of the same letter \(a\) with every intervening letter adjacent to \(a\). The intervening word may be empty.
\begin{lemma}\label{normal}
A word is shortest in \(W_T\) if and only if it is clean.
\end{lemma}
\begin{proof}
An offending pair can be brought together by commuting the first occurrence across the intervening letters; the pair then cancels, shortening the word by two.
For the converse, consider clean words modulo adjacent commuting swaps. These swaps preserve length. An occurrence of \(a\) is movable to the front if every letter before it commutes with \(a\). Define a toggle: delete that occurrence if it exists, and otherwise prepend \(a\). On clean words there is at most one such occurrence. The resulting word is clean: a newly offending pair after deletion would already have been offending, unless blocked by the deleted \(a\); but any letter of a pair straddling the deletion lies before that occurrence and commutes with \(a\), so it was not blocked by it. In the prepend case the only new offending pair would make \(a\) movable, contrary to the case assumption.

The toggle is well-defined on swap classes. Moving one commuting pair does not change the relative order of any noncommuting letters, so deletion of a movable occurrence gives the same swap class. A second toggle restores the original class: after deletion no other \(a\) can become movable, by cleanliness, and prepending the deleted letter recovers its original position by swaps. Toggles for adjacent distinct letters commute. Indeed, inserting or deleting one cannot block the other's movability, since they commute. Thus the four insertion/deletion cases yield identical classes after the two operations, with a possible swap of the inserted letters.

Consequently the toggles define involutive permutations satisfying the presentation relations, and hence an action of the actual presented group \(W_T\) on these classes. Each generator changes length by one. A clean word acting on the empty class gives its own class: read the product action from right to left, and prepend its successive letters; cleanliness precludes a deletion. If another word of length \(r\) represents the same group element, its action on the empty class reaches that same class, whose length is at most \(r\). The clean word is therefore shortest.
\end{proof}
In Lean these arguments are separated into \texttt{NormalForm.lean} and \texttt{TraceAction.lean}. Swap equivalence has no cancellation constructor. Length descends to its quotient, the toggles satisfy the Coxeter relations, and Mathlib's presentation universal property constructs the action. The group length is Mathlib's minimum generator-word length, not a language-based replacement definition. This formalization of standard graph-product normal-form reasoning supplies a semantic bridge and is not claimed as a new normal-form theorem.
\newpage
\section{Recognition and exact individual-word counts}
For a clean word \(w\), let \(S(w)\) be the letters that can be moved to the end using commuting swaps. Any two distinct letters in \(S(w)\) commute: in their original order the later occurrence must commute with the earlier one to move that earlier occurrence to the end. Thus \(S(w)\) is a clique.
Appending \(v\) fails cleanliness exactly when \(v\in S(w)\). Otherwise
\[
 S(wv)=\{v\}\cup(S(w)\cap N(v)).
\]
Indeed, an old letter remains movable exactly if it was movable before and commutes with the appended letter. Lemma~\ref{normal} now proves recognition of shortest words by the partial automaton
\[
 \delta(S,v)=\begin{cases}\text{failure},&v\in S,\\
 \{v\}\cup(S\cap N(v)),&v\notin S.
 \end{cases}
\]
The initial state is empty and every nonfailure state accepts. For a tree, its cliques are empty, single vertices, or edges; our two machines each have \(1+19+18=38\) states and 19 input letters. Failure may equivalently be a nonaccepting sink.

Here is a direct finite recurrence that makes the counting semantics explicit. For every accepting state \(S\), set \(b_0(S)=1\), and put
\[
 b_{k+1}(S)=\sum_{\substack{v\in V\\\delta(S,v)\ne\mathrm{failure}}}
 b_k(\delta(S,v)).
\]
Splitting a word into its first letter and its remaining \(k\) letters is a bijection onto this disjoint union of continuation sets. Induction proves that \(b_k(S)\) counts individual accepted input words of length \(k\) from \(S\). Summation is over letters, including multiplicities when distinct letters have the same target. Recognition then gives \(g_T(k)=b_k(\varnothing)\).

\texttt{ConcreteMachine.lean} contains all 722 transition entries for each graph, the state-meaning specification, and all 19 vectors \(b_0,\ldots,b_{18}\), each with 38 entries. Every transition and arithmetic recurrence is checked in the kernel. The two final empty-state entries are
\[
\begin{split}
 b_{18}^G(\varnothing)&=37769749324097905710910,\\
 b_{18}^H(\varnothing)&=37769749324097905710906.
\end{split}
\]
This proves Theorem~\ref{thm}. The equality of the earlier coefficients through 17 is recorded computationally in the independent review and is used only for the source-family comparison below; it is not a global optimality assertion.

\texttt{AutomatonSemantics.lean} proves recognition, including reversal to account for the implementation's reading convention. \texttt{WordCounting.lean} uses subtypes of functions \(\mathrm{Fin}(k)\to V\), establishes the split equivalence, and proves equality with actual shortest-word cardinality for every length. No quotient by commutation or group equality is used in the counted word type. \texttt{MachineInvariant.lean} identifies the successful state with the exact suffix state. The closed endpoint is \texttt{MathDB330264.conjecture\_false}.
\newpage
\section{Comparison with inspected constructions}
The mathematical separation is from the \emph{joint} three invariants, including the additional line-spectrum protection in [1, Conjecture 1.3]. It is one conjecture refutation. The following comparisons incorporate the independent significance review; they do not establish worldwide priority.

McKay [2, Lemma 4.1 and Theorem 4.2, pp.~223--225] proves bordered matrix identities and coalesces specified rooted 16-vertex limbs with a common rooted tree. Incidence identities transfer the construction to line spectra. Ciobanu--Kolpakov [1, Theorem 3.2 and Lemma 3.3] prove \emph{equal} geodesic growth for that source family. It cannot by that theorem alone give the unequal-growth conclusion here. McKay's Theorem 5.3 reconstructs a tree from a substantially larger collection of noncommutative adjacency/degree spectral data; it supplies no implication that the growth series distinguishes all nonisomorphic trees.

Godsil [3] is a slide deck dated 9 November 2007. Relevant labels are ``Spectrum of a 0-Sum'' (PDF p.~22), ``Corollary'' (p.~24), ``Another Pair'' (p.~27), and ``A Hint'' (p.~28), rather than invented theorem numbers. Ciobanu--Kolpakov [1, Theorem 5.1 and Lemma 5.2] use that rooted construction to obtain unequal growth but explicitly allow different line spectra. Their proof gives a nonzero discrepancy at length 10 for every nontrivial common rooted extension: its determining rooted counts are 8919523 and 8919522, multiplied by twice the positive attachment-root degree. The present pair agrees at length 10. That is an isomorphism-invariant obstruction to identifying this pair with that exact family, not an assertion that no composition of known constructions could refute the conjecture.

Godsil--McKay [4, Sections 2.1--2.7] give orthogonal switching and attachment-matrix constructions. Their proofs protect adjacency spectra, and in the specified switching setting complement spectra; they do not generally protect line spectra or establish the required geodesic inequality. We do not claim that the underlying graphs here are new.

The automaton is existing machinery [1, Definition 2.2; 5, Proposition 3.2; 7, Theorem 4.9]. Later equality formulas in [5, Theorem 3.7], [6, Theorems 1.1 and 2.2] and [7, Theorem 5.14] assume link-regularity. Our trees have degrees 1, 2 and 3, violating that condition already on singleton cliques. Those results do not directly supply this endpoint. The independent review read the indicated proofs, with the exception that [7, Theorem 6.4] was checked only at statement level and is not relied upon here.

The bounded searches and inspected primary bodies identified no complete prior theorem-to-refutation bridge. Search misses and the historical MathDB open label are not positive novelty evidence. Registry screening used the dated OpenAI/math head and complete preprints-subtree receipts described in the review; the root tree was truncated and screening was by paths, not all bodies. MathDB/Palomar access was denied and not retried. This manuscript inherits those limits rather than claiming fresh exhaustive clearance.
\newpage
\section{Source disposition, reviews and reproduction}
This is an AI-assisted research draft: manuscript preparation, proof engineering and dossier assembly used AI assistance. No human authorship or affiliation is asserted. The historical fidelity review passed source fidelity, mathematics and local proof/package gates with external-cache trust. The significance review gives a qualified significance pass, not worldwide novelty clearance or acceptance. Its journal-body condition is superseded by root's 4 November 2026 UTC source ruling, not by a claim that the journal body was obtained.

The authenticated arXiv versions are v1 (10 April 2015), v2 (6 March 2016), and v3 (17 January 2020). All retain Conjecture 1.3 unchanged in hypotheses and conclusion. The v3 abstract-page comment notes a typo correction in formula (3); the normalized v2/v3 body comparison concerns its adjugate numerator. The author-submitted v3 is positive post-publication retention evidence, not an authenticated journal manuscript. The 2016 version of record body remains unread. Root resolved the source-version hold as a qualified limitation on this basis. Bounded searches and access failures remain as documented; no exhaustive citation or registry clearance is claimed. No new network retrieval or author contact was made in this refresh.

\subsection*{Reproducible local certificate}
The current source-only packet is
\begin{quote}\small\path{frozen/330264-replay-fit-r2-source.zip}\end{quote}
Reused byte-identically; SHA256:
\begin{quote}\small\ttfamily
00c3713935d6a13f701cd3a26f9c7f1a866\newline
33e0eb60ff00a7ed0d100ea3d2fb5 (concatenate both lines).
\end{quote}
It contains 64 Lean modules. The former monolithic spectral proof is split into data, five-row certificates, identity assembly and topology modules. Changes are outside the mechanically computed statement closure: ten packet constants reached from the endpoint type by traversing types and bodies, excluding the endpoint proof as a seed. The closure was computed on the restructured compiled endpoint because the original module exceeded the wall cap. Its complete modules, \texttt{GroupSemantics} and \texttt{Endpoint}, remain byte-identical. Source-span comparisons and the preservation report document this boundary.

Six recursive definitions outside the closure receive \texttt{noncomputable} modifiers, with sections in four prerequisite modules, suppressing compiler-generated partial \texttt{\_unsafe\_rec} helpers. Their logical bodies are unchanged; removing precisely those annotations restores the original files. Nine matrix identities assemble checked row lemmas. The replay-fit-r2 report records fresh builds and explicit-module checker passes for all 64 modules, 874 local declarations without forbidden metadata, and endpoint axioms exactly \texttt{propext}, \texttt{Classical.choice}, \texttt{Quot.sound}. Its maximum reported call wall time is 52 seconds and sampled RSS 3,353,136 KiB, under two CPUs, 120 CPU/140 wall seconds and 6 GiB RSS. That lane discloses an earlier diagnostic-log write outside its boundary; this technical replay result is not an instruction-clean-execution claim.

\newpage
This manuscript lane performs safe extraction and byte checks, without running Lean, lake or leanchecker. See \texttt{REPRODUCE.md} and the preserved replay evidence. The original archive with hash prefix \texttt{2cab69a4} is retained only as historical provenance. The pinned toolchain is Lean 4.34.1, commit
\begin{quote}\small\ttfamily 5045d0056413266e57c625dcd7c365b10e377c52\end{quote}
Mathlib is pinned to
\begin{quote}\small\ttfamily d13f23b723b8a846827a245b89c10fc7d3f11612.\end{quote}

The historical dependency manifest pins 3,097 external source/olean pairs. Imported mathlib, package and core objects remain trusted cache inputs, not rebuilt proofs; matching hashes do not establish source-to-olean provenance. Implicit Init, runtime and plugin sidecars rely on the installed-toolchain assumption. This is an environment-pinned certificate, not a fully vendored portable build or full dependency rebuild.

Root's fresh capped replay and registry-refresh receipts are pending at this draft's evidence cutoff; no successful root replay or fresh registry clearance is inferred from partial logs. Root paper review and acceptance remain separate. Candidate credit is zero; the campaign context supplied for this refresh is 35/40 beyond 39. No publication, upload or shared ledger edit is authorized by this draft.
\newpage
\begin{thebibliography}{9}
\bibitem{ck} Laura Ciobanu and Alexander Kolpakov, \emph{Geodesic growth of right-angled Coxeter groups based on trees}. Inspected body: arXiv:1504.02774v3 (17 January 2020), Conjecture 1.3, Definitions 2.1--2.4, Theorem 3.2/Lemma 3.3 and Theorem 5.1/Lemma 5.2. Journal bibliographic record: J. Algebraic Combin. 44(2) (2016), 249--264, \href{https://doi.org/10.1007/s10801-016-0667-9}{doi:10.1007/s10801-016-0667-9}. Journal body not authenticated here. \url{https://arxiv.org/abs/1504.02774}
\bibitem{m} B. D. McKay, \emph{On the spectral characterisation of trees}, Ars Combinatoria 3 (1977), 219--232. \url{https://users.cecs.anu.edu.au/~bdm/papers/SpectralTrees.pdf}
\bibitem{g} C. Godsil, \emph{Are Almost All Graphs Cospectral?}, slide deck, 9 November 2007. \url{https://www.math.uwaterloo.ca/~cgodsil/pdfs/cospectral.pdf}
\bibitem{gm} C. D. Godsil and B. D. McKay, \emph{Constructing cospectral graphs}, Aequationes Math. 25 (1982), 257--268. \url{https://users.cecs.anu.edu.au/~bdm/papers/GodsilMcKayCospectral.pdf}
\bibitem{ac} Y. Antol\'in and L. Ciobanu, \emph{Geodesic growth in right-angled and even Coxeter groups}, European J. Combin. 34 (2013), 859--874; arXiv:1203.2752v2. \url{https://arxiv.org/abs/1203.2752}
\bibitem{af} Y. Antol\'in and I. Foniqi, \emph{Geodesic Growth of some 3-dimensional RACGs}, arXiv:2105.09751. \url{https://arxiv.org/abs/2105.09751}
\bibitem{ngp} L. Marjanski, V. Solon, F. Zheng and K. Zopff, \emph{Geodesic Growth of Numbered Graph Products}, Groups, Complexity, Cryptology 14(2) (2022), 2:1--2:32 (published 1 February 2023), DOI \href{https://doi.org/10.46298/jgcc.2022.14.2.10019}{10.46298/jgcc.2022.14.2.10019}. Publisher body: \url{https://gcc.episciences.org/10850/pdf}.
\end{thebibliography}
\end{document}
